\section{面向域内恢复训练的任务构建与执行验证}
\label{sec:domain-preparation}

\subsection{改进环节与验证目标}

前述端到端实验显示，记录选择最终会映射到完整来源页，多条不同选择路径可能形成相同的回答提示。围绕这一信息传递环节，本文进一步建立由程序直接提交条件化答案的任务：模型需要完成参数选择、跨记录对应、范围判断及结果组装，执行所得答案直接进入评价。模型动作因此与最终参数和条件保持明确联系，为域内恢复训练提供了可检验的目标。

本节报告该任务的数据构建、参考执行和训练接口验证结果。冻结材料包含576条合成训练参考、108条开发参考和四个新来源家族的24条评价参考。后续模型对照采用共享域内初始模型$I_{\mathrm{domain}}$，比较普通后缀续训$A_{\mathrm{domain}}$与恢复后缀续训$C_{\mathrm{domain}}$。初始模型由合成成功轨迹适配；两种续训从同一批训练失败中构造配对，并共享问题、数据顺序和监督预算。模型自由运行评价在计算环境恢复后开展。

\subsection{完整参考程序与类型核对}

全部708条参考程序在各自完整知识目录中执行通过，见表\ref{tab:offline-replays}。真实开发题覆盖12个完整查询目标，使用5--18次调用，包含跨表筛选、部件范围判断、配置对应及未知联合配置。每个目标均与原开发需求逐项比对；新来源24题另外完成43项来源量值不变量核对。合成参考的目标则由独立生成事实计算，672条结果全部一致。

\begin{table}[htbp]
\centering\small
\caption{冻结参考程序的CPU执行结果}
\label{tab:offline-replays}
\OfflineReplayTable
\end{table}

评分以完整结构化答案为主指标，同时记录成功提交、格式有效性、调用次数、无效调用和token消耗。表格行序和列的展示顺序不影响内容判断，选项和对应行可以重排，值与条件的配对关系保持不变。评分器统一已声明的兼容单位，检查范围端点、严格比较符和量值角色，并核对每列证据与冻结来源记录。空结果依据实际检索范围及条件判分；“资料未确立该组合”以独立判断状态计分；执行失败保留在问题分母中。

\subsection{恢复构造与上下文预算}

12道真实开发题的\OfflineRecoveryPairs{}组人工故障配对全部恢复至参考答案。两种故障分别验证成功动作占用新句柄与失败动作保持句柄不变的情况。监督检查同时覆盖共同前缀、分歧动作、工具观察和后缀目标，确认训练目标只落在指定助手位置。两侧监督token逐对匹配，最长测试上下文为\OfflineRecoveryContext{} token。

\begin{table}[htbp]
\centering\small
\caption{冻结参考轨迹的实际分词结果（token）}
\label{tab:offline-tokens}
\OfflineTokenTable
\end{table}

使用与基础模型对应的Qwen分词器检查全部708条参考轨迹，表\ref{tab:offline-tokens}列出各部分最大长度。最长参考上下文为\OfflineMaxContext{} token，依照运行前登记的规则选取统一\OfflineContextBudget{}上下文预算，每轮生成上限192 token、总生成上限2048 token、最多24次调用。训练与评价均保留完整问题和来源目录，超长或执行失败按固定规则记录。

这一阶段已形成可执行参考、评分器、分组隔离目录和恢复配对接口。接续实验先训练$I_{\mathrm{domain}}$并在开发部分检查组合查询能力，随后采集训练问题的真实失败，按配对数据开展两种续训，最后在固定24道基础题和12道组合补充题上分别报告完整答案及来源组表现。当前新增结果为来源复核与CPU验证；模型对照结果将在完成自由运行后纳入对应评价表。

\subsection{组合评价补充的执行与覆盖验证}
\label{sec:composition-validation}

在冻结的完整156条读法目录上，\CompositionQuestionCount{}道组合评价参考全部执行通过，合计\CompositionCalls{}次调用，每题5--11次。来源预期在程序编写前保存，核对覆盖完整量值、严格边界、工作模式、典型值限定及配置说明。对条件变化和必要信息移除进行的CPU检查进一步确认，程序结果由实际来源和查询条件共同决定。

\begin{table}[htbp]
\centering\small
\caption{基础与组合评价参考中的操作覆盖（题数）}
\label{tab:composition-operations}
\CompositionOperationTable
\end{table}

表\ref{tab:composition-operations}统计各题提交结果依赖路径中的操作。基础题主要支持新资料参数读取，组合题增加跨记录连接、显式筛选、范围判断及来源支持判断。一次操作可对应多个条件，同一题也可包含多种操作，因而各列的覆盖数分别统计。题目语义必要性由原件审阅确定，依赖路径检查用于确认相关操作实际参与答案生成。

实际分词检查得到最长完整上下文\CompositionMaxContext{} token、最大单步目标\CompositionMaxAction{} token、最大动作目标总量\CompositionMaxTargets{} token，全部符合既定16384上下文、192单步、2048总生成及24次调用预算。原24题、训练题和开发题的文件哈希保持一致。

最终域内模型评价采用两张独立结果表：原24题报告新来源基础应用表现，组合12题报告指定组合任务表现。各表保留全部问题，分别列出普通组与恢复组的正确数、改正与退步、完成率、来源分组和成本；同来源题的关联通过家族计数呈现。上述数据与执行验证为模型对照提供固定任务，后续按第一章实验矩阵开展实际训练及自由运行评价。

